% GENERATED FILE. Edit canonical lessons, then run npm run generate:books.
% canonical-source-hash: fnv1a64:bc0633c2438a339c
% canonical-lessons: PA-W12-greeting-choice, PA-W12-greeting-supported, PA-W12-closing-choice, PA-W12-closing-supported, PA-W12-boundary, PA-W12-two-sentence-order, PA-W12-two-sentence-supported, PA-W12-two-sentence-delayed, PA-W12-agreement-repair, PA-W12-spacing-repair, PA-W12-spelling-repair, PA-W12-punctuation-repair, PA-W12-wellbeing-supported, PA-W12-interests-supported, PA-W12-meeting-supported, PA-W12-body-blocks-delayed, PA-W12-message-supported, PA-W12-message-delayed, PA-W12-message-no-model

\chapter{Write to One Named Reader}
\label{ch:pa-a1-named-reader-message}
\hlchaptermodality{156}

\begin{chapteropening}
\textbf{By the end of this chapter:} \emph{I can choose and order already practised Gurmukhi chunks into an untimed independent 30–40-word message for one named reader, then check tone, meaning, agreement, spacing, spelling, and punctuation separately.}

The learner writes a no-model 30–40-word message to Manan for a stated purpose, independently selecting and ordering known chunks, then self-checks separate error dimensions without an exact-answer auto-grade.
\end{chapteropening}

\section[\texorpdfstring{\pa{ਨਮਸਤੇ} · \pa{ਸਤਿ} \pa{ਸ੍ਰੀ} … (namaste · sat …)}{ਨਮਸਤੇ · ਸਤਿ ਸ੍ਰੀ … (namaste · sat …)}]{A greeting that fits the reader}

\label{lesson:PA-W12-greeting-choice}



\begin{quote}
Read the two greetings you already know. Today you choose only the opening; the body of a message comes later.
\end{quote}

\begin{culture}
\textbf{\pa{ਸਤਿ} \pa{ਸ੍ਰੀ} \pa{ਅਕਾਲ}} is specifically a Sikh greeting. \textbf{\pa{ਨਮਸਤੇ}} is the broader greeting taught here. Do not guess a person's community from a Punjabi name. For a named reader whose Sikh context is not stated, choose \textbf{\pa{ਨਮਸਤੇ}}.
\end{culture}

\subsection*{Writing — choose before copying}
Imagine a note to Manan. Point to the broader opening, then copy it once in Gurmukhi: \textbf{\pa{ਨਮਸਤੇ}}. Keep the last vowel sign attached to its letter.

\begin{tcolorbox}[breakable,colback=teal!4,colframe=teal!35!black,title={Before you move on}]
One context choice is enough. The next tiny step places the reader's name.
\end{tcolorbox}
\section[\texorpdfstring{\pa{ਨਮਸਤੇ} \pa{ਮਨਨ।} (namaste manan)}{ਨਮਸਤੇ ਮਨਨ। (namaste manan)}]{Greet one named reader}

\label{lesson:PA-W12-greeting-supported}



\begin{quote}
Close the old six-field form and its value bank. For fictional card B · \pa{ਖ}, fill all six familiar labels from memory before opening the old page. Then compare one dimension at a time: selection, spelling, spaces, digits, and placement. Repair the first difference you find. This is a short recall of the form, not a new form lesson or a model for the message.

The name you retrieved is \textbf{\pa{ਮਨਨ}}. Say the greeting selected for a reader whose community is not specified. Now add that already-read name.
\end{quote}

\subsection*{Writing — a supported opening}
Copy \textbf{\pa{ਨਮਸਤੇ} \pa{ਮਨਨ।}} once. Leave one space between greeting and name, then put the Gurmukhi sentence mark immediately after the name. Point to those two boundaries; do not add a body sentence yet.

\subsection*{Your turn — change only the boundary}
Compare \textbf{\pa{ਨਮਸਤੇਮਨਨ।}} with \textbf{\pa{ਨਮਸਤੇ} \pa{ਮਨਨ।}} The first has no word boundary. Repair only that space. The greeting and name themselves do not change.

\begin{tcolorbox}[breakable,colback=teal!4,colframe=teal!35!black,title={Before you move on}]
This was a supported opening, not independent message writing.
\end{tcolorbox}
\section[\texorpdfstring{\pa{ਫਿਰ} \pa{ਮਿਲਾਂਗੇ} (phir milāṁge)}{ਫਿਰ ਮਿਲਾਂਗੇ (phir milāṁge)}]{End with the meeting you mean}

\label{lesson:PA-W12-closing-choice}



\begin{quote}
The opening is already chosen. Now choose only an ending for a note whose purpose is to meet Manan again.
\end{quote}

\begin{culture}
\textbf{\pa{ਫਿਰ} \pa{ਮਿਲਾਂਗੇ}} says “we will meet again.” It matches that purpose. The Sikh greeting \textbf{\pa{ਸਤਿ} \pa{ਸ੍ਰੀ} \pa{ਅਕਾਲ}} can also serve as a goodbye in its own context, but the reader's community has not been specified here.
\end{culture}

\subsection*{Writing — choose the ending}
Copy the everyday ending once: \textbf{\pa{ਫਿਰ} \pa{ਮਿਲਾਂਗੇ।}} Notice the word space and the mark after the final word. You are not assembling a message yet.

\begin{tcolorbox}[breakable,colback=teal!4,colframe=teal!35!black,title={Before you move on}]
Next you will place a courtesy word before this ending.
\end{tcolorbox}
\section[\texorpdfstring{\pa{ਧੰਨਵਾਦ।} \pa{ਫਿਰ} \pa{ਮਿਲਾਂਗੇ।} (dhanvād · phir …)}{ਧੰਨਵਾਦ। ਫਿਰ ਮਿਲਾਂਗੇ। (dhanvād · phir …)}]{Two small units at the end}

\label{lesson:PA-W12-closing-supported}



\begin{quote}
Read the everyday ending once. Then recall the old courtesy word \textbf{\pa{ਧੰਨਵਾਦ}}. They do different jobs: thanks first, parting second.
\end{quote}

\subsection*{Writing — supported closing}
Copy \textbf{\pa{ਧੰਨਵਾਦ।} \pa{ਫਿਰ} \pa{ਮਿਲਾਂਗੇ।}} as two short units. Put one space after the first sentence mark and one between \textbf{\pa{ਫਿਰ}} and \textbf{\pa{ਮਿਲਾਂਗੇ}}. No space goes before either mark.

\subsection*{Your turn — read the order}
Point to the word that thanks the reader. Point to the words that promise a later meeting. Keep this order; the next lesson will put an opening and this closing on separate lines.

\begin{tcolorbox}[breakable,colback=teal!4,colframe=teal!35!black,title={Before you move on}]
This is supported copying, not independent composition.
\end{tcolorbox}
\section[\texorpdfstring{\pa{ਨਮਸਤੇ} \pa{ਮਨਨ।} … (opening · closing)}{ਨਮਸਤੇ ਮਨਨ। … (opening · closing)}]{Leave room for the message}

\label{lesson:PA-W12-boundary}



\begin{quote}
Once more, close the old form bank and fill all six lines for fictional card B · \pa{ਖ} from memory. Reopen the old page only after the attempt; compare one dimension at a time and repair one difference. The form holds facts; do not paste its six labelled lines into the message.

The opening names the reader from that card. The closing thanks that reader and promises another meeting. Say which comes first before you write either.
\end{quote}

\subsection*{Writing — two ends, one empty middle}
Copy the two already practised ends. Leave the middle line empty for now:

\begin{quote}
\pa{ਨਮਸਤੇ} \pa{ਮਨਨ।}
\end{quote}

\begin{quote}
\_\_\_\_\_\_
\end{quote}

\begin{quote}
\pa{ਧੰਨਵਾਦ।} \pa{ਫਿਰ} \pa{ਮਿਲਾਂਗੇ।}
\end{quote}

Point to the blank. It is where facts for Manan will go, not another greeting.

\subsection*{Your turn — check the edges}
Check only the opening and closing. Does the opening name Manan? Does the closing put thanks before parting? Stop before filling the blank.

\begin{tcolorbox}[breakable,colback=teal!4,colframe=teal!35!black,title={Before you move on}]
A blank middle is a plan, not a finished message.
\end{tcolorbox}
\section[\texorpdfstring{\pa{ਮੇਰਾ} \pa{ਨਾਂ} \pa{ਅਮਨ} \pa{ਹੈ।} … (merā nāṁ aman …)}{ਮੇਰਾ ਨਾਂ ਅਮਨ ਹੈ। … (merā nāṁ aman …)}]{Name first, language second}

\label{lesson:PA-W12-two-sentence-order}



\begin{quote}
The old connected reading gave you these two facts about Aman. This lesson changes only their order for a reader meeting Aman on the page.
\end{quote}

\begin{tcolorbox}[breakable,skin=enhanced,colback=orange!6,colframe=orange!45!black,boxrule=0.5pt,arc=1mm,left=6pt,right=6pt,top=4pt,bottom=4pt,fonttitle=\bfseries,title={Read this}]
\begin{quote}
\pa{ਮੈਂ} \pa{ਪੰਜਾਬੀ} \pa{ਬੋਲਦਾ} \pa{ਹਾਂ।}
\end{quote}

\begin{quote}
\pa{ਮੇਰਾ} \pa{ਨਾਂ} \pa{ਅਮਨ} \pa{ਹੈ।}
\end{quote}

Both sentences are already taught. Point to the name sentence, then to the language sentence. The reader needs to know who \textbf{\pa{ਮੈਂ}} is.
\end{tcolorbox}

\subsection*{Writing — order before composing}
Copy \textbf{\pa{ਮੇਰਾ} \pa{ਨਾਂ} \pa{ਅਮਨ} \pa{ਹੈ।}} first and \textbf{\pa{ਮੈਂ} \pa{ਪੰਜਾਬੀ} \pa{ਬੋਲਦਾ} \pa{ਹਾਂ।}} second. Do not join them or change their words. One sentence boundary is the whole task.

\begin{tcolorbox}[breakable,colback=teal!4,colframe=teal!35!black,title={Before you move on}]
You ordered two old chunks; no new grammar was needed.
\end{tcolorbox}
\section[\texorpdfstring{\pa{ਮੇਰਾ} \pa{ਨਾਂ} \pa{ਅਮਨ} \pa{ਹੈ।} \pa{ਮੈਂ} … (two supported …)}{ਮੇਰਾ ਨਾਂ ਅਮਨ ਹੈ। ਮੈਂ … (two supported …)}]{Copy one small body block}

\label{lesson:PA-W12-two-sentence-supported}



\begin{quote}
Read the name sentence and the language sentence in the order chosen last time. Their meanings are already familiar.
\end{quote}

\subsection*{Writing — supported two-sentence block}
Copy this block, not the whole message:

\begin{quote}
\pa{ਮੇਰਾ} \pa{ਨਾਂ} \pa{ਅਮਨ} \pa{ਹੈ।} \pa{ਮੈਂ} \pa{ਪੰਜਾਬੀ} \pa{ਬੋਲਦਾ} \pa{ਹਾਂ।}
\end{quote}

Check \textbf{\pa{ਹੈ}} at the end of the name sentence and \textbf{\pa{ਹਾਂ}} at the end of the speaker sentence. Keep one space after the first sentence mark.

\subsection*{Your turn — one contrast}
Point to \textbf{\pa{ਹੈ}}, then \textbf{\pa{ਹਾਂ}}. Aman is named in the first sentence; the speaker says “I” in the second. Do not repair a different dimension yet.

\begin{tcolorbox}[breakable,colback=teal!4,colframe=teal!35!black,title={Before you move on}]
The model stays visible today. Delayed recall comes next.
\end{tcolorbox}
\section[(delayed name and language pair)]{Recall the same two sentences later}

\label{lesson:PA-W12-two-sentence-delayed}



\begin{quote}
Close the page with the copied block. Take a short pause. The only cues here are “Aman's name” and “the language Aman speaks.”
\end{quote}

\subsection*{Writing — delayed two-sentence block}
From memory, write the name sentence first and the speaker-language sentence second. Leave a clear space after the first sentence mark. Do not reopen the model until both sentences are on the page.

\subsection*{Your turn — compare after writing}
Now reopen the supported lesson. Check the sentence order, then \textbf{\pa{ਹੈ}} and \textbf{\pa{ਹਾਂ}}, then the mark and following space. Repair one difference at a time.

\begin{tcolorbox}[breakable,colback=teal!4,colframe=teal!35!black,title={Before you move on}]
This is delayed recall of two known sentences, not the final message.
\end{tcolorbox}
\section[\texorpdfstring{\pa{ਹੈ} · \pa{ਹਾਂ} (hai · hāṁ)}{ਹੈ · ਹਾਂ (hai · hāṁ)}]{Repair only the endings}

\label{lesson:PA-W12-agreement-repair}



\begin{quote}
Recall two old endings: the name sentence says \textbf{\pa{ਹੈ}}; the speaker sentence with \textbf{\pa{ਮੈਂ}} says \textbf{\pa{ਹਾਂ}}. Today no word order or spelling changes.
\end{quote}

\begin{tcolorbox}[breakable,skin=enhanced,colback=orange!6,colframe=orange!45!black,boxrule=0.5pt,arc=1mm,left=6pt,right=6pt,top=4pt,bottom=4pt,fonttitle=\bfseries,title={Read this}]
\begin{quote}
\pa{ਮੇਰਾ} \pa{ਨਾਂ} \pa{ਅਮਨ} \pa{ਹਾਂ।} \pa{ਮੈਂ} \pa{ਪੰਜਾਬੀ} \pa{ਬੋਲਦਾ} \pa{ਹੈ।}
\end{quote}

Both endings have traded places. The names, spaces, sentence marks, and other words are already right. Point only to the two endings.
\end{tcolorbox}

\subsection*{Writing — repair the endings}
Rewrite the two sentences, changing \textbf{\pa{ਹਾਂ}} after \textbf{\pa{ਅਮਨ}} to \textbf{\pa{ਹੈ}} and \textbf{\pa{ਹੈ}} after \textbf{\pa{ਬੋਲਦਾ}} to \textbf{\pa{ਹਾਂ}}. Leave every other part untouched.

\begin{tcolorbox}[breakable,colback=teal!4,colframe=teal!35!black,title={Before you move on}]
The next pass checks spaces, not endings.
\end{tcolorbox}
\section[\texorpdfstring{\pa{ਨਮਸਤੇ} \pa{ਮਨਨ।} · \pa{ਮੇਰਾ} \pa{ਨਾਂ} (two visible word …)}{ਨਮਸਤੇ ਮਨਨ। · ਮੇਰਾ ਨਾਂ (two visible word …)}]{Give each word its space}

\label{lesson:PA-W12-spacing-repair}



\begin{quote}
The greeting and name were copied separately earlier. Today only the spaces are wrong; do not change a letter or sentence mark.
\end{quote}

\begin{tcolorbox}[breakable,skin=enhanced,colback=orange!6,colframe=orange!45!black,boxrule=0.5pt,arc=1mm,left=6pt,right=6pt,top=4pt,bottom=4pt,fonttitle=\bfseries,title={Read this}]
\begin{quote}
\pa{ਨਮਸਤੇਮਨਨ।} \pa{ਮੇਰਾਨਾਂ} \pa{ਅਮਨ} \pa{ਹੈ।}
\end{quote}

Point after the greeting and after \textbf{\pa{ਮੇਰਾ}}. Those are two word boundaries, not two spelling errors.
\end{tcolorbox}

\subsection*{Writing — repair only spaces}
Insert one space in each crowded place. Keep the sentence marks attached to the preceding words. Compare with the old supported greeting only afterward.

\begin{tcolorbox}[breakable,colback=teal!4,colframe=teal!35!black,title={Before you move on}]
The next pass looks inside one word.
\end{tcolorbox}
\section[\texorpdfstring{\pa{ਪੰਜਾਬੀ} (panjābī)}{ਪੰਜਾਬੀ (panjābī)}]{One missing mark inside a known word}

\label{lesson:PA-W12-spelling-repair}



\begin{quote}
Say the familiar word for Punjabi, then picture its Gurmukhi spelling. This is a one-word check, not a new sentence.
\end{quote}

\begin{tcolorbox}[breakable,skin=enhanced,colback=orange!6,colframe=orange!45!black,boxrule=0.5pt,arc=1mm,left=6pt,right=6pt,top=4pt,bottom=4pt,fonttitle=\bfseries,title={Read this}]
In \textbf{\pa{ਮੈਂ} \pa{ਪਜਾਬੀ} \pa{ਬੋਲਦਾ} \pa{ਹਾਂ।}}, the spaces, verb, and final mark are right. The word for the language has lost its nasal mark. Compare the taught \textbf{\pa{ਪੰਜਾਬੀ}} with \textbf{\pa{ਪਜਾਬੀ}} and point to the missing mark.
\end{tcolorbox}

\subsection*{Writing — repair one spelling}
Write only \textbf{\pa{ਪੰਜਾਬੀ}}. Check the mark above \textbf{\pa{ਪ}}. Do not copy the rest of the sentence or turn this into an agreement pass.

\begin{tcolorbox}[breakable,colback=teal!4,colframe=teal!35!black,title={Before you move on}]
The next pass checks where sentences end.
\end{tcolorbox}
\section[\texorpdfstring{\pa{।} (sentence boundary)}{। (sentence boundary)}]{Put the stops back}

\label{lesson:PA-W12-punctuation-repair}



\begin{quote}
You have already checked endings, word spaces, and one spelling. Today those are all fixed. Look only for where each familiar sentence stops.
\end{quote}

\begin{tcolorbox}[breakable,skin=enhanced,colback=orange!6,colframe=orange!45!black,boxrule=0.5pt,arc=1mm,left=6pt,right=6pt,top=4pt,bottom=4pt,fonttitle=\bfseries,title={Read this}]
\begin{quote}
\pa{ਨਮਸਤੇ} \pa{ਮਨਨ} \pa{ਮੇਰਾ} \pa{ਨਾਂ} \pa{ਅਮਨ} \pa{ਹੈ} \pa{ਮੈਂ} \pa{ਪੰਜਾਬੀ} \pa{ਬੋਲਦਾ} \pa{ਹਾਂ}
\end{quote}

The greeting ends after the reader's name. The name fact ends after \textbf{\pa{ਹੈ}}. The language fact ends after \textbf{\pa{ਹਾਂ}}. No word needs to move.
\end{tcolorbox}

\subsection*{Writing — repair only punctuation}
Insert \textbf{\pa{।}} at those three stops. Keep every word and existing word space. Leave one space after each new mark that is followed by another sentence.

\begin{tcolorbox}[breakable,colback=teal!4,colframe=teal!35!black,title={Before you move on}]
Four separate repair passes are complete; the message can now grow in supported blocks.
\end{tcolorbox}
\section[\texorpdfstring{\pa{ਤੁਸੀਂ} \pa{ਕਿਵੇਂ} \pa{ਹੋ}? · \pa{ਮੈਂ} … (tusī̃ kivē̃ hō …)}{ਤੁਸੀਂ ਕਿਵੇਂ ਹੋ? · ਮੈਂ … (tusī̃ kivē̃ hō …)}]{Ask the reader, then share one small answer}

\label{lesson:PA-W12-wellbeing-supported}



\begin{quote}
The opening already names Manan. The old respectful question asks how the reader is. The old \textbf{\pa{ਮੈਂ}} sentence gives Aman's own wellbeing.
\end{quote}

\begin{tcolorbox}[breakable,skin=enhanced,colback=orange!6,colframe=orange!45!black,boxrule=0.5pt,arc=1mm,left=6pt,right=6pt,top=4pt,bottom=4pt,fonttitle=\bfseries,title={Read this}]
\begin{quote}
\pa{ਤੁਸੀਂ} \pa{ਕਿਵੇਂ} \pa{ਹੋ}?
\end{quote}

\begin{quote}
\pa{ਮੈਂ} \pa{ਠੀਕ}-\pa{ਠਾਕ} \pa{ਹਾਂ।}
\end{quote}

The first line faces Manan; the second speaks for Aman. The question mark belongs to the first line, the Gurmukhi sentence mark to the second.
\end{tcolorbox}

\subsection*{Writing — supported wellbeing block}
Copy the two lines once. Keep \textbf{\pa{ਹੋ}} with \textbf{\pa{ਤੁਸੀਂ}} and \textbf{\pa{ਹਾਂ}} with \textbf{\pa{ਮੈਂ}}. No new reply or verb form is needed.

\begin{tcolorbox}[breakable,colback=teal!4,colframe=teal!35!black,title={Before you move on}]
This is one supported pair, not a full message.
\end{tcolorbox}
\section[\texorpdfstring{\pa{ਮੈਨੂੰ} \pa{ਪੰਜਾਬੀ} \pa{ਪਸੰਦ} \pa{ਹੈ।} … (mainū̃ panjābī …)}{ਮੈਨੂੰ ਪੰਜਾਬੀ ਪਸੰਦ ਹੈ। … (mainū̃ panjābī …)}]{Two interests, one familiar frame}

\label{lesson:PA-W12-interests-supported}



\begin{quote}
The old liking frame begins \textbf{\pa{ਮੈਨੂੰ}} and ends \textbf{\pa{ਪਸੰਦ} \pa{ਹੈ}}. Name the two already taught things Aman likes: Punjabi and reading.
\end{quote}

\begin{tcolorbox}[breakable,skin=enhanced,colback=orange!6,colframe=orange!45!black,boxrule=0.5pt,arc=1mm,left=6pt,right=6pt,top=4pt,bottom=4pt,fonttitle=\bfseries,title={Read this}]
\begin{quote}
\pa{ਮੈਨੂੰ} \pa{ਪੰਜਾਬੀ} \pa{ਪਸੰਦ} \pa{ਹੈ।}
\end{quote}

\begin{quote}
\pa{ਮੈਨੂੰ} \pa{ਪੜ੍ਹਨਾ} \pa{ਪਸੰਦ} \pa{ਹੈ।}
\end{quote}

The frame does not change. The first middle is a language; the second is reading. \textbf{\pa{ਪੜ੍ਹਨਾ}} was taught earlier, including its joined spelling.
\end{tcolorbox}

\subsection*{Writing — supported interest block}
Copy both short sentences. Point to \textbf{\pa{ਮੈਨੂੰ}} and \textbf{\pa{ਪਸੰਦ} \pa{ਹੈ}} in each, then check the one word that changed. Do not invent a new interest for this pass.

\begin{tcolorbox}[breakable,colback=teal!4,colframe=teal!35!black,title={Before you move on}]
The model remains visible today.
\end{tcolorbox}
\section[\texorpdfstring{\pa{ਅਸੀਂ} \pa{ਜਲਦੀ} \pa{ਮਿਲਾਂਗੇ।} (asī̃ jaldī milāṁge)}{ਅਸੀਂ ਜਲਦੀ ਮਿਲਾਂਗੇ। (asī̃ jaldī milāṁge)}]{A plan before the goodbye}

\label{lesson:PA-W12-meeting-supported}



\begin{quote}
The old reading said that Aman and Manan will meet soon. The closing says they will meet again. This lesson writes only the plan, not both sentences.
\end{quote}

\subsection*{Writing — supported meeting line}
Copy \textbf{\pa{ਅਸੀਂ} \pa{ਜਲਦੀ} \pa{ਮਿਲਾਂਗੇ।}} once. Point to \textbf{\pa{ਅਸੀਂ}} for the two people and \textbf{\pa{ਜਲਦੀ}} for when. Keep the sentence mark after the plan.

\subsection*{Your turn — place the line}
In a note to Manan, this line goes after Aman's facts and before \textbf{\pa{ਧੰਨਵਾਦ।} \pa{ਫਿਰ} \pa{ਮਿਲਾਂਗੇ।}} Do not write the whole note yet; point to where the line belongs.

\begin{tcolorbox}[breakable,colback=teal!4,colframe=teal!35!black,title={Before you move on}]
All body pieces are now visible separately. Delay comes before assembly.
\end{tcolorbox}
\section[(delayed message body blocks)]{Retrieve the body in small blocks}

\label{lesson:PA-W12-body-blocks-delayed}



\begin{quote}
Close the four supported pages. Take a short pause. Keep three English cues: who Aman is, how the two people are and what Aman likes, and when they meet.
\end{quote}

\subsection*{Writing — delayed blocks, not a whole letter}
Write the old name-and-language pair. Pause. Write the old reader question and Aman's wellbeing reply. Pause. Write the two old liking sentences and the one meeting plan. Leave the opening and closing for later. Do not reopen the supported pages until every block has an attempt.

\subsection*{Your turn — compare one block at a time}
Reopen one earlier lesson for each block. Check order and meaning first, then endings, word spaces, spelling, and punctuation in separate passes. Repair only the first difference found on a pass.

\begin{tcolorbox}[breakable,colback=teal!4,colframe=teal!35!black,title={Before you move on}]
This is delayed retrieval of old blocks. It does not yet score an independently assembled message.
\end{tcolorbox}
\section[\texorpdfstring{\pa{ਨਮਸਤੇ} \pa{ਮਨਨ।} · \pa{ਫਿਰ} … (supported …)}{ਨਮਸਤੇ ਮਨਨ। · ਫਿਰ … (supported …)}]{One visible message from old blocks}

\label{lesson:PA-W12-message-supported}



\begin{quote}
Every line below has already been practised alone or in a short pair. Today the only new demand is assembly for Manan. Read the opening and closing first, then find the small body blocks between them.
\end{quote}

\begin{tcolorbox}[breakable,skin=enhanced,colback=orange!6,colframe=orange!45!black,boxrule=0.5pt,arc=1mm,left=6pt,right=6pt,top=4pt,bottom=4pt,fonttitle=\bfseries,title={Read this}]
\begin{quote}
\pa{ਨਮਸਤੇ} \pa{ਮਨਨ।} \pa{ਤੁਸੀਂ} \pa{ਕਿਵੇਂ} \pa{ਹੋ}?
\end{quote}

\begin{quote}
\pa{ਮੇਰਾ} \pa{ਨਾਂ} \pa{ਅਮਨ} \pa{ਹੈ।} \pa{ਮੈਂ} \pa{ਪੰਜਾਬੀ} \pa{ਬੋਲਦਾ} \pa{ਹਾਂ।} \pa{ਮੈਂ} \pa{ਠੀਕ}-\pa{ਠਾਕ} \pa{ਹਾਂ।}
\end{quote}

\begin{quote}
\pa{ਮੈਨੂੰ} \pa{ਪੰਜਾਬੀ} \pa{ਪਸੰਦ} \pa{ਹੈ।} \pa{ਮੈਨੂੰ} \pa{ਪੜ੍ਹਨਾ} \pa{ਪਸੰਦ} \pa{ਹੈ।}
\end{quote}

\begin{quote}
\pa{ਅਸੀਂ} \pa{ਜਲਦੀ} \pa{ਮਿਲਾਂਗੇ।} \pa{ਧੰਨਵਾਦ।} \pa{ਫਿਰ} \pa{ਮਿਲਾਂਗੇ।}
\end{quote}

Count the space-separated words: \textbf{30}. The greeting and question face the reader; the middle shares familiar facts and interests; the meeting plan leads into thanks and goodbye.
\end{tcolorbox}

\subsection*{Writing — supported assembly}
Copy one line at a time, pausing after each line. Keep the reader's question mark distinct from the Gurmukhi sentence marks. Check meaning and order before spelling. This is a visible model, not independent composition.

\begin{tcolorbox}[breakable,colback=teal!4,colframe=teal!35!black,title={Before you move on}]
The model disappears for the delayed attempt.
\end{tcolorbox}
\section[(delayed named-reader message)]{Recall the model only after a pause}

\label{lesson:PA-W12-message-delayed}



\begin{quote}
Close the supported message. Wait briefly. Keep only the broad English memory cues: reader greeting, wellbeing question, name and language, interests, meeting plan, thanks and goodbye.
\end{quote}

\subsection*{Writing — delayed assembly}
Recall the earlier Gurmukhi message in those small blocks. Pause after the name-and-language block and again after the interests. Finish the whole attempt before reopening the model. Do not treat this recall of a model as independent choice of content.

\subsection*{Your turn — compare in passes}
Now reopen the supported page. Check that the note addresses Manan and ends with the meeting plan, then check sentence order, endings, word spaces, spelling, and punctuation one at a time. Count the words last.

\begin{tcolorbox}[breakable,colback=teal!4,colframe=teal!35!black,title={Before you move on}]
The next lesson changes the task: it gives a reader and purpose, not a message to reproduce.
\end{tcolorbox}
\section[(independent named-reader message)]{Manan writes a new message to Aman}

\label{lesson:PA-W12-message-no-model}



\begin{quote}
Close all earlier pages. You already know the short pieces needed for a note. This time there is no sentence bank or sample to reproduce.
\end{quote}

\subsection*{Writing — independent named-reader message}
As Manan, write Aman a \textbf{30–40-word} Gurmukhi message. Greet Aman and ask how he is. Identify yourself as Manan. Include the already taught liking for water and at least one other familiar interest. Put your own wellbeing reply \emph{after} the interest sentence or sentences. Suggest meeting again and close. Choose any additional familiar fact you need to reach the word band. Do not reopen a page or use a dictionary, translator, or spell-checker before finishing your draft.

There is no exam clock. Write one short block, pause, and continue. The earlier supported model went from Aman to Manan, did not include the water liking, and put the wellbeing reply earlier. Recalling it unchanged cannot meet this new task.

\subsection*{Your turn — self-check after the draft}
Count space-separated words and keep the result between 30 and 40. Check one dimension at a time: the named reader and appropriate tone; meaning and sentence order; agreement; spacing; Gurmukhi spelling; then punctuation. Only after recording those checks may you reopen earlier pages to repair a word. A different well-formed message can be valid; no single exact string is the answer to free composition.

\begin{tcolorbox}[breakable,colback=teal!4,colframe=teal!35!black,title={Before you move on}]
Record the word count and one repair you made. This untimed checkpoint does not award timed-paper credit or whole-A1 readiness; two complete mocks and human validation remain separate work.
\end{tcolorbox}
